\section{Preliminaries}
\label{sec:prelim}

\subsection{Conventions}

We identify $\R^2$ with $\C$, writing $x=(x_1,x_2)=x_1+ix_2$ when convenient,
and $\bar x$ for the complex conjugate; $|x|$ is the Euclidean norm; for a matrix or an array of numbers, $|\cdot|$ is the
Euclidean norm of the array of its entries, and for a matrix $M$ we write
$\|M\|_{\mathrm{op}}$ for its operator norm. For $r>0$, $B_r=\{|x|<r\}$ and
$\overline B_r$ is its closure; $|S|$ is the Lebesgue measure of a measurable
set $S\subset\R^2$. The overline thus denotes complex conjugation, of numbers and
of functions, and the closure of a set; it is also part of the names of some fixed
objects, such as the radial background $\bomega$ of Section~\ref{sec:ring} and the
constants $\Gone$, $\sbar$ and $\bar T$. Which meaning is intended will always be
clear from the context. We use polar
coordinates $x=r(\cos\theta,\sin\theta)$, with $\er=(\cos\theta,\sin\theta)$
and $\etheta=(-\sin\theta,\cos\theta)$. For $\theta\in\R$, $\Orot_\theta$ is
the rotation of $\R^2$ by the angle $\theta$ (multiplication by $e^{i\theta}$),
and $\Orot=\Orot_{2\pi/3}$. For an integer $m\ge1$, a function $g$ on $\R^2$ is
\textit{$m$-fold rotationally symmetric} if and only if $g\circ\Orot_{2\pi/m}=g$ almost
everywhere, and \textit{radial} if and only if $g\circ\Orot_\theta=g$ almost
everywhere for all $\theta$. If $\theta\mapsto g(re^{i\theta})$ is integrable
on $[0,2\pi]$ for almost every $r>0$ (by Fubini's theorem in polar
coordinates, this holds for every $g\in L^1_{\mathrm{loc}}(\R^2)$) and
$\ell\in\Z$, the \textbf{angular mode} $\ell$ of $g$ is
$g_\ell(r)=\frac1{2\pi}\int_0^{2\pi}g(re^{i\theta})e^{-i\ell\theta}\dd\theta$,
defined for almost every $r>0$. If $g\in L^2(\R^2)$, then
$g=\sum_{\ell\in\Z}g_\ell(r)e^{i\ell\theta}$ in $L^2(\R^2)$, and $g$ is
$m$-fold rotationally symmetric if and only if $g_\ell=0$ for all $\ell\notin m\Z$.

The Biot-Savart kernel is $\KBS(x)=x^\perp/(2\pi|x|^2)$ with
$x^\perp=(-x_2,x_1)$, and $\nabla^\perp=(-\partial_2,\partial_1)$. For a
vorticity $\omega$ we write $v=\KBS\ast\omega$ whenever the integral converges.
If $\omega\in C_c^\infty(\R^2)$, then the \textbf{stream function}
$\psi(x)=\frac1{2\pi}\int\log|x-y|\,\omega(y)\dd y$ satisfies
$v=\nabla^\perp\psi$, $\Delta\psi=\omega$ and $\partial_1v_2-\partial_2v_1=\omega$
(Lemma~\ref{lem:BS-basic}(d) below). In polar coordinates these sign
conventions read $v\cdot\er=-r^{-1}\partial_\theta\psi$ and
$v\cdot\etheta=\partial_r\psi$. For a vector field $u$ we write $\nabla u$ for
the matrix $(\partial_lu_i)_{i,l=1,2}$, so that $\operatorname{tr}\nabla
u=\divop u$.

Functions are complex valued unless we say otherwise; vorticities and
velocities are real. For a multi-index $\gamma=(\gamma_1,\gamma_2)\in\N_0^2$ we
write $\partial^\gamma=\partial_1^{\gamma_1}\partial_2^{\gamma_2}$,
$|\gamma|=\gamma_1+\gamma_2$, $\zeta^{2\gamma}=\zeta_1^{2\gamma_1}\zeta_2^{2\gamma_2}$
for $\zeta\in\R^2$,
$\beta\le\gamma$ if $\beta_1\le\gamma_1$ and $\beta_2\le\gamma_2$, and
$\binom\gamma\beta=\binom{\gamma_1}{\beta_1}\binom{\gamma_2}{\beta_2}$; here
$\N=\{1,2,\dots\}$, $\N_0=\N\cup\{0\}$ and $m\N=\{m,2m,\dots\}$. For an integer
$k\ge0$ and an open set $V\subset\R^2$, $H^k(V)$ is the Sobolev space with norm
$\|g\|_{H^k(V)}^2=\sum_{|\gamma|\le k}\|\partial^\gamma g\|_{L^2(V)}^2$, and
$H^k=H^k(\R^2)$. Sobolev norms of vector and matrix fields are summed over the
entries; for instance $\|\nabla u\|_{H^k}^2=\sum_{i,l}\|\partial_lu_i\|_{H^k}^2$.
For a vector or matrix field $F$ on $V$ and $1\le q\le\infty$ we put
$\|F\|_{L^q(V)}=\|\,|F|\,\|_{L^q(V)}$, where $|F|$ is the Euclidean norm of the
array of entries; for $q=2$ this agrees with the sum over the entries. For a
vector field $F$ on $V$ we put
$\|F\|_{W^{1,\infty}(V)}=\operatorname{ess\,sup}_V|F|+\operatorname{ess\,sup}_V|\nabla F|$.

For a function or vector field $h$ of class $C^q$ near a point $x$, we write
$\nabla^qh(x)$ for the array of all partial derivatives of order $q$ of the
components of $h$ at $x$, one entry for each ordered $q$-tuple of directions,
so that $|\nabla^qh(x)|$ is its Euclidean norm. For an integer $k\ge0$, a closed
set $U\subset\R^2$ and $h$ of class $C^k$ on a neighborhood of $U$, we put
$\|h\|_{C^k(U)}=\sum_{q=0}^k\sup_{x\in U}|\nabla^qh(x)|$;
in particular $|\partial^\gamma h(x)|\le\|h\|_{C^k(U)}$ for $x\in U$ and
$|\gamma|\le k$, and $C^k(U)$ denotes the space of such $h$ with finite norm.
The Fourier transform is $\hat g(\zeta)=\int g(x)e^{-ix\cdot\zeta}\dd x$, so that
$\|\partial^\gamma g\|_{L^2(\R^2)}^2=(2\pi)^{-2}\int\zeta^{2\gamma}|\hat g(\zeta)|^2\dd\zeta$
by Plancherel's theorem.

We put $\jap{t}=(1+t^2)^{1/2}$ for $t\in\R$ and $\log_+x=\max\{\log x,0\}$ for
$x>0$. For $x>0$ and an integer $p\ge0$ we put
\begin{equation}
\label{eq:Ip}
\Ip{p}{x}=\int_0^\infty e^{-xu}(1+u)^p\dd u,
\end{equation}
a finite (as $x>0$), decreasing function of $x$.

Constants that enter a later choice of parameters carry their own names, and
their dependence is stated where they are defined. The letter $C$ denotes a
constant that enters no later choice; its dependence is stated each time, and
within one proof its value may change from line to line. For a function $u$
and a nonnegative function $g$, $u=O(g)$ as $r\to r_0$ means that $|u|\le Cg$ in
a neighborhood of $r_0$ for some constant $C$.

Let $I$ be an interval and $X$ a Banach space. A function $F\colon I\to X$ is
\textit{strongly measurable} if it is an almost everywhere limit of simple
functions $\sum_{i=1}^k\mathbf 1_{E_i}(t)x_i$, with $E_i\subset I$ measurable
and $x_i\in X$, and \textit{weakly measurable} if $t\mapsto\Lambda(F(t))$ is
measurable for every $\Lambda\in X^*$. We use Pettis's theorem \cite{Pettis}:
if $X$ is separable, then $F$ is strongly measurable if and only if it is weakly
measurable. The Bochner spaces $L^\infty(I;X)$ and $L^1(I;X)$ consist of the
strongly measurable $F$ for which $t\mapsto\|F(t)\|_X$ is essentially bounded,
respectively integrable, modulo equality almost everywhere.

Let $\Omega\subset\R^2$ be a measurable set. For
$b\in[1,\infty)$, $L^\infty(I;L^b(\Omega))$ and $L^1(I;L^b(\Omega))$ are the
Bochner spaces, whose elements we identify with measurable functions on
$I\times\Omega$; a measurable $u$ on $I\times\Omega$ lies in them if and only if
$t\mapsto\|u(t)\|_{L^b(\Omega)}$ is essentially bounded, respectively
integrable, by Pettis's theorem: $L^b(\Omega)$ is separable, and for
$\phi\in L^{b'}(\Omega)$, $1/b+1/b'=1$, the function $t\mapsto\int u(t)\phi$,
defined for almost every $t$, is measurable by Tonelli's theorem applied to the
positive and negative parts of the real and imaginary parts of $u\phi$. We write $L^\infty(I;L^\infty(\Omega))$
for $L^\infty(I\times\Omega)$, and $L^1(I;L^\infty(\Omega))$ for the set of
measurable $u$ on $I\times\Omega$ with $\int_I\|u(t)\|_{L^\infty(\Omega)}\dd t<\infty$.
Intersections such as $L^\infty(I;(L^1\cap L^\infty)(\Omega))$ are intersections of these
spaces; a space with $W^{1,a}(\Omega)$ in place of $L^b(\Omega)$ consists of the
$w$ for which $w$ and its distributional gradient in $x$ lie in the
corresponding space with $L^a(\Omega)$; the subscript $\mathrm{loc}$ on the time
interval means that the condition holds on $J\times\Omega$ for every compact
interval $J\subset I$, on $\R^2$ that it holds on $I\times B_R$ for every $R>0$,
and on both that it holds on $J\times B_R$ for all such $J$ and $R$.

For $T>0$ and $1\le q\le\infty$, the space $X_q(T)$ of Section~1, written
$L^\infty_t(L^1\cap L^q)_x$ there, is the Bochner space
$L^\infty\big(0,T;(L^1\cap L^q)(\R^2)\big)$ just described, with
$\|\omega\|_{(L^1\cap L^q)(\R^2)}=\|\omega\|_{L^1(\R^2)}+\|\omega\|_{L^q(\R^2)}$; and
$\DT=\{\varphi\in C_c^\infty(\R\times\R^2):\supp\varphi\subset(-\infty,T)\times\R^2\}$.
$\LipO$ is the set of Lipschitz functions $\Ren\colon\R\to\R$ with $\Ren(0)=0$.
For $g\in L^1(\R^2)$ we put $\Circ_g(r)=\int_{B_r}g$.

A $C_0$-semigroup $(S(t))_{t\ge0}$ on a Banach space $X$ is a family of bounded
operators with $S(0)$ the identity, $S(t+s)=S(t)S(s)$, and $t\mapsto S(t)x$ continuous for
all $x\in X$; its generator $A$ is $Ax=\lim_{h\downarrow0}h^{-1}(S(h)x-x)$ on the
set $D(A)$ of $x$ for which the limit exists.

\subsection{The Biot-Savart Law}

We split $\KBS=\Knear+\Kfar$ with $\Knear=\mathbf 1_{B_1}\KBS$. By
Lemma~\ref{lem:BS-basic}(a) below,
$\|\Knear\|_{L^s(\R^2)}^s=(2\pi)^{1-s}/(2-s)$ for $s\in[1,2)$, and
$|\Kfar(z)|\le\min\{1,|z|^{-1}\}/(2\pi)$ for $z\in\R^2$.

\begin{lemma}
\label{lem:BS-basic}
The kernel $\KBS$ has the following properties.
\begin{enumerate}[label=\textup{(\alph*)}]
\item For $z\ne0$, $i,l\in\{1,2\}$ and $\rho>0$, $|\KBS(z)|=(2\pi|z|)^{-1}$,
  $|\partial_l(\KBS)_i(z)|\le(2\pi|z|^2)^{-1}$ and
  $\int_{B_\rho}|\KBS|\dd z=\rho$.
\item Let $g\in (L^1\cap L^\infty)(\R^2)$. Then for every $x\in\R^2$,
  \[
  \big(|\KBS|\ast|g|\big)(x)\le\Big(\frac{2\|g\|_{L^1(\R^2)}\|g\|_{L^\infty(\R^2)}}{\pi}\Big)^{1/2} .
  \]
  In particular the integral defining $\KBS\ast g(x)$ converges absolutely for
  every $x$, and
  $\|\KBS\ast g\|_{L^\infty(\R^2)}\le\big(2\|g\|_{L^1(\R^2)}\|g\|_{L^\infty(\R^2)}/\pi\big)^{1/2}$.
\item Let $\rho,R>0$ and let $g\in L^2(\R^2)$ vanish outside $B_\rho$. Then
  $\KBS\ast g$ is defined almost everywhere, and
  $\|\KBS\ast g\|_{L^2(B_R)}\le(R+\rho)\|g\|_{L^2(\R^2)}$.
\item Let $g\in C_c^\infty(\R^2)$ and $u=\KBS\ast g$. Then $u\in C^\infty(\R^2)$,
  $\partial^\gamma u=\KBS\ast\partial^\gamma g$ for every $\gamma$,
  $u=\nabla^\perp\psi$ with $\psi(x)=\frac1{2\pi}\int\log|x-y|\,g(y)\dd y$,
  $\Delta\psi=g$, $\divop u=0$ and $\partial_1u_2-\partial_2u_1=g$. Moreover
  $\sum_{i,l=1}^2\|\partial_l\partial^\gamma u_i\|_{L^2(\R^2)}^2=\|\partial^\gamma g\|_{L^2(\R^2)}^2$
  for every $\gamma$, and hence $\|\nabla u\|_{H^k}=\|g\|_{H^k}$ for every
  integer $k\ge0$.
\item Let $g\in L^1(\R^2)$ and $R>0$. Then
  $\int_{B_R}\big(|\KBS|\ast|g|\big)\dd x\le R\,\|g\|_{L^1(\R^2)}$. In particular the
  integral defining $\KBS\ast g$ converges absolutely almost everywhere,
  $\KBS\ast g\in L^1_{\mathrm{loc}}(\R^2)$, and $\divop(\KBS\ast g)=0$ in
  $\mathcal D'(\R^2)$.
\end{enumerate}
\end{lemma}

\begin{proof}
(a) The first identity holds because $|x^\perp|=|x|$, and integration in polar
coordinates gives $\int_{B_\rho}(2\pi|z|)^{-1}\dd z=\rho$. Since
$2\pi\KBS(z)=(-z_2,z_1)/|z|^2$, the derivatives $2\pi\partial_l(\KBS)_i$ are,
up to sign, the numbers
$\partial_l(z_k/|z|^2)=\delta_{lk}|z|^{-2}-2z_lz_k|z|^{-4}$, whose modulus is
$|z_{k'}^2-z_k^2|/|z|^4\le|z|^{-2}$ if $l=k$ (where $k'\ne k$) and
$2|z_1z_2|/|z|^4\le|z|^{-2}$ if $l\ne k$.

(b) We may assume $g\ne0$. For arbitrary $r>0$, (a) gives, after splitting the
integral into the near field $\{|z|<r\}$ and the far field $\{|z|\ge r\}$,
\[
\begin{aligned}
2\pi\big(|\KBS|\ast|g|\big)(x)&\le\|g\|_{L^\infty(\R^2)}\int_{B_r}\frac{\dd z}{|z|}
+\frac1r\int_{|z|\ge r}|g(x-z)|\dd z\\
&\le2\pi r\|g\|_{L^\infty(\R^2)}+\frac{\|g\|_{L^1(\R^2)}}{r},
\end{aligned}
\]
and choosing $r=(\|g\|_{L^1(\R^2)}/(2\pi\|g\|_{L^\infty(\R^2)}))^{1/2}$ gives
\[
2\pi(|\KBS|\ast|g|)(x)\le2(2\pi\|g\|_{L^1(\R^2)}\|g\|_{L^\infty(\R^2)})^{1/2},
\]
which is the claim.
The rest follows from $|\KBS\ast g|\le|\KBS|\ast|g|$.

(c) For $|x|<R$ and $|x-z|<\rho$, we have $|z|<\rho+R$, so, as $g$ vanishes
outside $B_\rho$, $\KBS\ast g(x)=\int_{B_{R+\rho}}\KBS(z)g(x-z)\dd z$ for $x\in B_R$, and Minkowski's
integral inequality and (a) give
$\big\|\int_{B_{R+\rho}}|\KBS(z)||g(\cdot-z)|\dd z\big\|_{L^2(B_R)}\le(R+\rho)\|g\|_{L^2(\R^2)}$.
In particular the integral converges absolutely almost everywhere on $B_R$, for
every $R>0$.

(d) Let $\supp g\subset B_\rho$. Since $\KBS$ and $\log|\cdot|$ are locally
integrable, differentiation under the integral sign shows that $u,\psi\in
C^\infty$, $\partial^\gamma u=\KBS\ast\partial^\gamma g$ and
$\partial^\gamma\psi=\frac1{2\pi}\log|\cdot|\ast\partial^\gamma g$. Writing
$\partial_jg(x-z)=-\partial_{z_j}[g(x-z)]$ and integrating by parts on
$\{|z|>\varepsilon\}$, with a boundary term at most
$\varepsilon|\log\varepsilon|\sup|g|\to0$, we get
$\partial_j\psi(x)=\int\frac{z_j}{2\pi|z|^2}g(x-z)\dd z$. Hence
$u=\nabla^\perp\psi$, so $\divop u=-\partial_1\partial_2\psi+\partial_2\partial_1\psi=0$
and $\Delta\psi=\partial_1u_2-\partial_2u_1$. Fix $x$, put
$h(z)=g(x-z)$, and use polar coordinates $z=re^{i\theta}$, in which
$z\cdot\nabla_z=r\partial_r$.
Since $\partial_lu=\KBS\ast\partial_lg$ and $(\partial_lg)(x-z)=-\partial_{z_l}h(z)$,
\[
(\partial_1u_2-\partial_2u_1)(x)=-\int\frac{z\cdot\nabla_zh}{2\pi|z|^2}\dd z
=-\frac1{2\pi}\int_0^{2\pi}\!\!\int_0^\infty\partial_rh\dd r\dd\theta=g(x),
\]
because $(\KBS)_2\partial_1h-(\KBS)_1\partial_2h=z\cdot\nabla_zh/(2\pi|z|^2)$ and
$\int_0^\infty\partial_rh\dd r=-h(0)=-g(x)$; the integrand is integrable, as
$|\partial_rh|\le|\nabla h|$ and $h$ has compact support. Hence $\Delta\psi=g$.

For $|x|\ge2\rho$ and $|y|<\rho$ we have $|x-y|\ge|x|/2$; differentiating under
the integral sign, with $|\nabla\log|z||=|z|^{-1}$ and
$|\partial_i\partial_j\log|z||=|\partial_j(z_i/|z|^2)|\le|z|^{-2}$ (proof of
(a)), gives $|\nabla\psi(x)|\le\|g\|_{L^1(\R^2)}/(\pi|x|)$ and
$|\partial_i\partial_j\psi(x)|\le2\|g\|_{L^1(\R^2)}/(\pi|x|^2)$ for $|x|\ge2\rho$.
Two integrations by parts on $B_r$
give $\int_{B_r}\sum_{i,j}(\partial_i\partial_j\psi)^2\dd x
=\int_{B_r}(\Delta\psi)^2\dd x+\mathcal B_r$, where $\mathcal B_r$ is a sum of
integrals over $\partial B_r$ of products of a first and a second derivative of
$\psi$; for $r\ge2\rho$ the bounds above give
$|\mathcal B_r|\le Cr\cdot r^{-1}r^{-2}\|g\|_{L^1(\R^2)}^2$ with $C$ absolute. Letting
$r\to\infty$ gives $\sum_{i,j}\|\partial_i\partial_j\psi\|_{L^2(\R^2)}^2=\|g\|_{L^2(\R^2)}^2$.
The entries of $\nabla u$ are, up to sign, the second derivatives of $\psi$, so
$\sum_{i,l}\|\partial_lu_i\|_{L^2(\R^2)}^2=\|g\|_{L^2(\R^2)}^2$. Replacing $g$ by
$\partial^\gamma g$, whose velocity is $\partial^\gamma u$, and summing over
$|\gamma|\le k$ gives the rest. For complex $g$ we apply all of this to the
real and imaginary parts.

(e) Let $D$ be a disk of radius $R$. Then $|D\setminus B_R|=|B_R\setminus D|$.
Since $|\KBS(z)|=(2\pi|z|)^{-1}$ by (a), $|\KBS|\le(2\pi R)^{-1}$ on $D\setminus B_R$
and $|\KBS|>(2\pi R)^{-1}$ on $B_R\setminus D$; so
$\int_D|\KBS|\le\int_{B_R}|\KBS|=R$, by (a). For fixed $y$, the substitution
$z=x-y$ maps $B_R$ onto the disk $D_y=\{z:|z+y|<R\}$ of radius $R$, so Tonelli's
theorem gives
$\int_{B_R}(|\KBS|\ast|g|)\dd x=\int|g(y)|\int_{D_y}|\KBS(z)|\dd z\dd y\le R\|g\|_{L^1(\R^2)}$.
Hence $|\KBS|\ast|g|$ is finite almost everywhere, and
$\KBS\ast g\in L^1_{\mathrm{loc}}(\R^2)$ because $|\KBS\ast g|\le|\KBS|\ast|g|$.
Take $g_n\in C_c^\infty(\R^2)$ with $g_n\to g$ in $L^1(\R^2)$. This bound, applied to
$g_n-g$, gives $\KBS\ast g_n\to\KBS\ast g$ in $L^1(B_R)$ for every $R>0$, hence
in $\mathcal D'(\R^2)$; since $\divop(\KBS\ast g_n)=0$ by (d),
$\divop(\KBS\ast g)=0$ in $\mathcal D'(\R^2)$.
\end{proof}

\begin{lemma}
\label{lem:velocity}
Let $4/3<p<\infty$, $p'=p/(p-1)$, and $R>0$. There is $C=C(p,R)$ such that
the following hold.
\begin{enumerate}[label=\textup{(\alph*)}]
\item For $\omega\in (L^1\cap L^p)(\R^2)$ the integral defining $\KBS\ast\omega$
  converges absolutely almost everywhere and, with $v=\KBS\ast\omega$,
  \[
  \|v\|_{L^{p'}(B_R)}+\|v\|_{L^2(B_R)}\le C\|\omega\|_{(L^1\cap L^p)(\R^2)},\qquad
  \int_{\R^2}|\omega|\,\big(|\KBS|\ast|\omega|\big)\dd x\le C\|\omega\|_{(L^1\cap L^p)(\R^2)}^2 ;
  \]
  the constant in the second inequality can be chosen depending only on $p$,
  independently of $R$.
\item Let $\omega_n$ be real and $\omega_n\rightharpoonup\omega$ weakly in $L^p(\R^2)$, with
  $\sup_n\|\omega_n\|_{L^1(\R^2)}\le M$. Then $\|\omega\|_{L^1(\R^2)}\le M$ and
  $\KBS\ast\omega_n\to\KBS\ast\omega$ strongly in $L^{p'}(B_R)$.
\item If $\omega\in (L^1\cap L^2)(\R^2)$, then $v=\KBS\ast\omega$ satisfies $\nabla v\in L^2(\R^2)$ and
  $\|\nabla v\|_{L^2(\R^2)}=\|\omega\|_{L^2(\R^2)}$.
\end{enumerate}
\end{lemma}

\begin{proof}
(a) Put $r(p)=p'$ and $s(p)=p/(2p-2)$ if $p\le2$, and $r(p)=p$ and $s(p)=1$ if
$p>2$. Since $p>4/3$, we have $s(p)\in[1,2)$, so $\Knear\in L^{s(p)}(\R^2)$. Moreover
$1+1/r(p)=1/s(p)+1/p$, so Young's inequality gives
$\||\Knear|\ast|\omega|\|_{L^{r(p)}(\R^2)}\le\|\Knear\|_{L^{s(p)}(\R^2)}\|\omega\|_{L^p(\R^2)}$.
Also $\||\Kfar|\ast|\omega|\|_{L^\infty(\R^2)}\le\|\omega\|_{L^1(\R^2)}/(2\pi)$. Hence
$|\KBS|\ast|\omega|$ is finite almost everywhere. Since
$r(p)\ge\max\{2,p'\}$, H\"older's inequality on $B_R$ gives the velocity bound.
For the last bound, $r(p)'\in[1,p]$, so
$\|\omega\|_{L^{r(p)'}(\R^2)}\le\|\omega\|_{(L^1\cap L^p)(\R^2)}$ by interpolation, and
$\int|\omega|(|\KBS|\ast|\omega|)\le\|\omega\|_{L^{r(p)'}(\R^2)}\||\Knear|\ast|\omega|\|_{L^{r(p)}(\R^2)}
+\|\omega\|_{L^1(\R^2)}^2/(2\pi)$.

(b) For $L>0$ and $g\in L^\infty(B_L)$, the function $g\mathbf 1_{B_L}$ lies in
$L^{p'}(\R^2)$, so $|\int_{B_L}\omega g|=\lim_n|\int_{B_L}\omega_ng|\le
M\|g\|_{L^\infty(\R^2)}$. The limit $\omega$ is real, as a weak limit of real functions.
Taking $g=\sgn\omega$ (with $\sgn s=s/|s|$ for $s\ne0$ and $\sgn0=0$) and
letting $L\to\infty$ gives $\|\omega\|_{L^1(\R^2)}\le M$.

By hypothesis and the bound just proved, $\omega_n$ and $\omega$ lie in
$(L^1\cap L^p)(\R^2)$, with $L^1(\R^2)$ norms at most $M$. Put
$N=\sup_n\|\omega_n\|_{L^p(\R^2)}+\|\omega\|_{L^p(\R^2)}<\infty$ (weakly convergent
sequences are bounded), let $0<\varepsilon<1<L$, and split
$\KBS=K^{(1)}+K^{(2)}+K^{(3)}$ with $K^{(1)}=\KBS\mathbf 1_{B_\varepsilon}$,
$K^{(2)}=\KBS\mathbf 1_{B_L\setminus B_\varepsilon}$ and
$K^{(3)}=\KBS\mathbf 1_{\R^2\setminus B_L}$. By Young's inequality, as in (a), and
H\"older's inequality on $B_R$,
$\|K^{(1)}\ast(\omega_n-\omega)\|_{L^{p'}(B_R)}\le|B_R|^{1/p'-1/r(p)}\|K^{(1)}\|_{L^{s(p)}(\R^2)}N$,
and $\|K^{(1)}\|_{L^{s(p)}(\R^2)}\to0$ as $\varepsilon\to0$ because $s(p)<2$. Next,
$|K^{(3)}\ast(\omega_n-\omega)|\le2M/(2\pi L)$ everywhere, since
$\|\omega_n\|_{L^1(\R^2)},\|\omega\|_{L^1(\R^2)}\le M$; so its $L^{p'}(B_R)$ norm is at most
$|B_R|^{1/p'}M/(\pi L)$, uniformly in $n$. Finally, for
$x\in B_R$ the function $K^{(2)}(x-\cdot)$ is bounded by $(2\pi\varepsilon)^{-1}$,
since $|\KBS(z)|=(2\pi|z|)^{-1}$ and $|z|\ge\varepsilon$ on the support of
$K^{(2)}$, and vanishes outside $\{y:|x-y|<L\}$; hence it lies in $L^{p'}(\R^2)$, so $K^{(2)}\ast\omega_n(x)\to K^{(2)}\ast\omega(x)$, with
$|K^{(2)}\ast\omega_n(x)|\le\|K^{(2)}\|_{L^{p'}(\R^2)}N$; dominated convergence gives
$K^{(2)}\ast\omega_n\to K^{(2)}\ast\omega$ in $L^{p'}(B_R)$. Letting first
$n\to\infty$ and then $\varepsilon\to0$ and $L\to\infty$ gives (b).

(c) Choose $\omega_j\in C_c^\infty(\R^2)$ with $\omega_j\to\omega$ in $L^1(\R^2)$ and in
$L^2(\R^2)$, and put $v_j=\KBS\ast\omega_j$. By Lemma~\ref{lem:BS-basic}(d), applied
to $\omega_j-\omega_i$ and to $\omega_j$, $\|\nabla v_j-\nabla v_i\|_{L^2(\R^2)}=\|\omega_j-\omega_i\|_{L^2(\R^2)}$
and $\|\nabla v_j\|_{L^2(\R^2)}=\|\omega_j\|_{L^2(\R^2)}$, so $\nabla v_j\to F$ in $L^2(\R^2)$ with
$\|F\|_{L^2(\R^2)}=\lim_j\|\omega_j\|_{L^2(\R^2)}=\|\omega\|_{L^2(\R^2)}$. By
Lemma~\ref{lem:BS-basic}(e), $v_j\to v$ in $L^1_{\mathrm{loc}}(\R^2)$, so
$\nabla v_j\to\nabla v$ in $\mathcal D'$, and $\nabla v=F$.
\end{proof}

\begin{definition}
\label{def:multipole}
Let $m\ge1$ be an integer and let $g\in L^1(\R^2;\C)$ have compact support in
$\R^2\setminus\{0\}$. Writing $y\in\R^2$ as a complex number, we put
\[
\pout_m[g]=\frac1{2\pi}\int g(y)\,(\bar y)^m\dd y,\quad
\pin_m[g]=\frac1{2\pi}\int g(y)\,y^{-m}\dd y,\quad
\Qm[g]=\int g(y)\,y^{m}\dd y .
\]
\end{definition}

These are the multipole moments of $g$. In polar coordinates $y=se^{i\vartheta}$,
since $(\bar y)^m=s^me^{-im\vartheta}$, $y^{-m}=s^{-m}e^{-im\vartheta}$ and
$y^m=s^me^{im\vartheta}$, they read, with the angular modes $g_\ell$ of
Section~\ref{sec:prelim},
\[
\pout_m[g]=\int_0^\infty g_m(s)\,s^{m+1}\dd s,\qquad
\pin_m[g]=\int_0^\infty g_m(s)\,s^{1-m}\dd s,\qquad
\Qm[g]=2\pi\int_0^\infty g_{-m}(s)\,s^{m+1}\dd s ;
\]
so they see only the angular modes $\pm m$ of $g$. Let $g$ be real with
$\supp g\subset\{\rho_1\le|y|\le\rho_2\}$, $0<\rho_1\le\rho_2$, and let
$\psi(x)=\frac1{2\pi}\int\log|x-y|\,g(y)\dd y$ be its stream function. Expanding
$\log|x-y|=\log|x|+\Re\log(1-y/x)$ for $|x|>\rho_2$, and
$\log|x-y|=\log|y|+\Re\log(1-x/y)$ for $|x|<\rho_1$, in powers of $y/x$ and
$x/y$ gives
\[
\begin{aligned}
\psi(x)&=\frac{\int g}{2\pi}\log|x|-\Re\sum_{m\ge1}\frac{\pout_m[g]}m(\bar x)^{-m}&&(|x|>\rho_2),\\
\psi(x)&=\frac1{2\pi}\int g(y)\log|y|\dd y-\Re\sum_{m\ge1}\frac{\pin_m[g]}m\,x^m&&(|x|<\rho_1),
\end{aligned}
\]
since $\overline{\Qm[g]}=2\pi\pout_m[g]$ for real $g$. Thus $\pout_m[g]$ and
$\pin_m[g]$ give the field of $g$ in the angular modes $\pm m$ outside and inside its
support.

\subsection{Weak Solutions}

\begin{definition}
\label{def:equimeasurable}
Two real functions $g,h\in L^1(\R^2)$ are \textit{equimeasurable} if and only if
$|\{g>s\}|=|\{h>s\}|$ and $|\{g<-s\}|=|\{h<-s\}|$ for every $s>0$.
\end{definition}

\begin{definition}
\label{def:weak}
Let $4/3<q<\infty$, $T>0$ and $\omega^{\mathrm{in}}\in (L^1\cap L^q)(\R^2)$. A
\textit{weak solution with datum $\omega^{\mathrm{in}}$} in $X_q(T)$ is a function
$\omega\in X_q(T)$ such that, with $v(t)=\KBS\ast\omega(t)$,
\begin{equation}
\label{eq:weak}
\int_0^T\!\!\int\omega\,(\partial_t\varphi+v\cdot\nabla\varphi)\dd x\dd t
+\int\omega^{\mathrm{in}}\varphi(0,\cdot)\dd x=0
\qquad\text{for all }\varphi\in\DT .
\end{equation}
It is \textit{renormalized} if and only if, in addition, for every
$\Ren\in C^1(\R)$ with $\Ren'$ bounded and $\Ren(0)=0$,
\begin{equation}
\label{eq:renorm}
\int_0^T\!\!\int\Ren(\omega)\,(\partial_t\varphi+v\cdot\nabla\varphi)\dd x\dd t
+\int\Ren(\omega^{\mathrm{in}})\varphi(0,\cdot)\dd x=0
\qquad\text{for all }\varphi\in\DT .
\end{equation}
For $T_0\in(0,\infty]$, a (renormalized) weak solution \textit{on $[0,T_0)$} is a
function on $[0,T_0)$ whose restriction to $[0,T)$ is a (renormalized) weak
solution in $X_q(T)$ for every $T<T_0$.
\end{definition}

The integrals in Equations~\eqref{eq:weak} and~\eqref{eq:renorm} converge
absolutely, because $|\Ren(\omega)|\le\|\Ren'\|_{L^\infty(\R)}|\omega|$ and
$|\omega v|\le|\omega|(|\KBS|\ast|\omega|)$, and Lemma~\ref{lem:velocity}(a),
with $p=q$ and $R=1$, gives
$\|\omega(t)v(t)\|_{L^1(\R^2)}\le C(q,1)\operatorname{ess\,sup}_{0<s<T}\|\omega(s)\|_{(L^1\cap L^q)(\R^2)}^2$
for almost every $t\in(0,T)$.
